% =============================================================================
% Section 1 — Introduction
% =============================================================================
\section{Introduction}
\label{sec:introduction}

Anyone who has attended a rock concert, stood near a pipe organ in full
voice, or visited an electronic dance music (EDM) event will recognise
the sensation: a deep, visceral awareness of the bass that seems to
bypass the ears and arrive directly in the gut.  Music fans describe
this as ``feeling the bass,'' and sound engineers deliberately design
sub-woofer arrays to maximise it~\cite{Leventhall2007, Moller2004}.
But what, physically, is being felt?

Two plausible pathways exist.  First, \emph{airborne acoustic coupling}:
sound pressure waves impinge on the body wall, deform it, and excite
internal resonances.  Second, \emph{structural (whole-body) vibration}:
sound transmitted through the floor, seats, and skeleton shakes the
torso directly.  For infrasound below \SI{20}{\hertz}, our companion
study~\cite{Mace2026browntone} showed that airborne coupling to the
abdominal cavity is negligibly weak: the Helmholtz number
$ka \approx 0.011$ at the $n = 2$ flexural resonance
(${\sim}\SI{4}{\hertz}$) imposes a $(ka)^n$ penalty that reduces
effective driving pressure by four orders of magnitude.  Whole-body
vibration, by contrast, couples directly at full amplitude, explaining
the well-documented gastrointestinal disturbances reported under ISO~2631
exposure conditions~\cite{iso2631, Griffin1990}.

The sub-bass band (\SIrange{20}{80}{\hertz}), however, has not been
examined in this framework.  This frequency range is of particular
interest for three reasons:
%
\begin{enumerate}
  \item \textbf{Concert relevance.}  The dominant spectral content of
    modern amplified music---kick drums, bass guitars, synthesised
    sub-bass---lies in the \SIrange{30}{80}{\hertz}
    band~\cite{Gunderson2016, Opperman2006}.

  \item \textbf{Coupling transition.}  At \SI{20}{\hertz},
    $ka \approx 0.06$; at \SI{80}{\hertz}, $ka \approx 0.23$.  The
    acoustic coupling coefficient $(ka)^n$ is therefore
    \numrange{30}{400} times larger than at the \SI{4}{\hertz} flexural
    resonance.  It is not \emph{a priori} obvious that the coupling
    remains negligible.

  \item \textbf{Psychoacoustic bridge.}  Sub-bass frequencies straddle
    the boundary between hearing and whole-body vibration
    perception~\cite{Todd2000, iso226}.  Understanding whether airborne
    sound can produce visceral sensation at these frequencies connects
    structural acoustics to psychoacoustics in a way that has not been
    attempted quantitatively.
\end{enumerate}

Pipe organs provide an extreme case.  The 32-foot rank of a large organ
produces a fundamental at \SI{16}{\hertz} with harmonics extending
through the sub-bass band at sound pressure levels approaching
\SI{105}{\decibel} in the nave~\cite{Bickerdike1966}.  Congregants have
long reported physical sensations during full-organ passages, and these
reports predate electrical amplification by centuries.

In the electronic music domain, purpose-built sub-woofer arrays routinely
produce \SIrange{110}{120}{\decibel} in the \SIrange{30}{60}{\hertz}
band~\cite{Opperman2006}.  The ``drop''---a sudden onset of heavy
sub-bass after a quiet passage---is a central aesthetic element of EDM,
and its visceral impact is a deliberate design goal.

Despite widespread anecdotal evidence and significant cultural
importance, no study has quantitatively assessed whether airborne
sub-bass at concert levels can produce physiologically meaningful tissue
displacement in the human abdomen.  Existing work on whole-body vibration
perception~\cite{Griffin1990, Mansfield2005} focuses on the mechanical
pathway (floor and seat vibration), while studies of concert sound
exposure~\cite{Gunderson2016, Opperman2006} focus on hearing damage
rather than visceral perception.

The present study addresses this gap.  We extend the fluid-filled
viscoelastic oblate spheroidal shell model of Mace and
Mace~\cite{Mace2026browntone} to the \SIrange{20}{80}{\hertz} band,
compute energy-consistent tissue displacements at literature-derived
concert sound pressure levels using the Junger \& Feit reciprocity
formulation, and compare these with both whole-body vibration perception
thresholds from ISO~2631-1~\cite{iso2631} and estimates of
structure-borne (floor/seat) vibration.  Rather than asking simply
``does airborne coupling work?''\ we quantitatively partition the
visceral sensation into its airborne and structure-borne components,
providing a complete pathway budget for the ``feel it in your gut''
experience.
